\documentclass[11pt,a4paper,landscape]{article}

\usepackage{fontawesome5}
\providecommand{\faPencilAlt}{\faIcon{pencil-alt}}

\usepackage{styles/couleurs}
\usepackage{styles/mise_en_page}
\usepackage{styles/boites}
\usepackage{styles/maths}
\usepackage{styles/figures}

\usepackage{enumitem}
\usepackage{array}
\usepackage{tabularx}
\usepackage{tikz}
\usepackage{graphicx}
\usepackage{amssymb}
\usepackage{multicol}

\usetikzlibrary{calc}

\geometry{
  a4paper,
  landscape,
  top=0.8cm,
  bottom=0.8cm,
  left=1cm,
  right=1cm
}

\pagestyle{empty}

\setlength{\parindent}{0pt}
\setlength{\parskip}{0.12em}
\setlength{\columnsep}{0.85cm}

\setlist[enumerate]{
  leftmargin=*,
  itemsep=0.08em,
  topsep=0.08em,
  parsep=0pt
}

\titlespacing*{\section}
  {0pt}
  {0.35em}
  {0.12em}

\newcommand{\soustitreexo}[1]{%
  {\sffamily\bfseries\color{TitreC}#1}\par
  \vspace{0.04em}
}

\begin{document}

\begin{center}
  {\sffamily\LARGE\bfseries\color{TitreC}
  Correction -- Séance d'AP : Statistiques}\\[-0.05em]

  {\sffamily\normalsize\bfseries\color{SavoirC}
  Vocabulaire, effectifs, moyenne et fréquences}
\end{center}

\vspace{-0.45em}

\begin{multicols*}{2}

%==================================================
\section*{1. Vocabulaire et effectifs}
%==================================================

\soustitreexo{Exercice 1}

On étudie les livres lus par \(24\) élèves.

\begin{enumerate}[label=\alph*)]
  \item \textbf{Population :} les \(24\) élèves interrogés.
  \item \textbf{Caractère étudié :} le nombre de livres lus pendant les vacances.
  \item \textbf{Valeurs prises :} \(0;\ 1;\ 2;\ 3;\ 4;\ 5\).
  \item La valeur \(2\) apparaît \(6\) fois : son \textbf{effectif est \(6\)}.
\end{enumerate}

\soustitreexo{Exercice 2}

\begin{enumerate}[label=\alph*)]
  \item Effectif des élèves venant en bus : \(\boxed{11}\).
  \item Effectif total :
  \[
  6+11+4+7=\boxed{28}.
  \]
  \item Élèves ne venant pas en bus :
  \[
  28-11=\boxed{17}.
  \]
\end{enumerate}

%==================================================
\section*{2. Lire un document}
%==================================================

\soustitreexo{Exercice 3}

\begin{enumerate}[label=\alph*)]
  \item Effectif en natation : \(\boxed{9}\).
  \item L'activité d'effectif le plus faible est la \textbf{danse}, avec \(\boxed{5}\) élèves.
  \item Effectif total :
  \[
  8+5+9+6=\boxed{28}.
  \]
\end{enumerate}

%==================================================
\section*{3. Moyenne}
%==================================================

\soustitreexo{Exercice 4 -- Une petite série}

\[
M=\frac{9+5+5+9+9+9+5+6+2+6}{10}
=\frac{65}{10}
=\boxed{6{,}5}.
\]

\textbf{Interprétation :} la moyenne des notes de Charles est de
\(\boxed{6{,}5/20}\).

\soustitreexo{Exercice 5 -- Une série plus longue}

La somme des \(31\) températures est \(424\).

\[
M=\frac{424}{31}\approx\boxed{13{,}7}.
\]

\textbf{Interprétation :} durant ce mois, la température quotidienne
a été en moyenne d'environ \(\boxed{13{,}7\,^\circ\mathrm{C}}\).

\columnbreak

%==================================================
\section*{3. Moyenne -- suite}
%==================================================

\soustitreexo{Exercice 6 -- Des données regroupées}

Chaque pointure doit être comptée autant de fois que l'indique son effectif :

\[
M=
\frac{
35\times7+
36\times9+
37\times5+
38\times2+
39\times7
}{
7+9+5+2+7
}.
\]

\[
M=\frac{1103}{30}\approx\boxed{36{,}8}.
\]

\textbf{Interprétation :} la pointure moyenne des membres du club
est d'environ \(\boxed{36{,}8}\).

\remarque{Une moyenne n'est pas obligatoirement une valeur observée :
aucun membre du club n'a besoin de chausser du \(36{,}8\).}


\end{multicols*}

\end{document}
